% ============================================================
\chapter{Introduction \`a la Th\'eorie Ergodique}
\label{ch:ergodique}
% ============================================================

\section{Introduction}

La th\'eorie ergodique\index{th\'eorie ergodique} \'etudie le comportement
statistique des syst\`emes dynamiques \`a long terme. Elle r\'epond \`a la
question fondamentale : les moyennes temporelles co\"incident-elles avec les
moyennes spatiales ? Ce chapitre pr\'esente les transformations pr\'eservant
la mesure, les th\'eor\`emes ergodiques fondamentaux et leurs applications.

\begin{intuition}
Consid\'erons un gaz dans une bo\^ite. Plut\^ot que de suivre chaque
mol\'ecule, on mesure des grandeurs macroscopiques (temp\'erature, pression)
qui sont des moyennes. L'hypoth\`ese ergodique de Boltzmann affirme que la
moyenne temporelle d'une observable le long d'une trajectoire \'egale sa
moyenne sur tout l'espace des phases. C'est le fondement de la m\'ecanique
statistique.
\end{intuition}

\section{Transformations pr\'eservant la mesure}

\begin{definition}[Syst\`eme dynamique mesur\'e]
\index{syst\`eme dynamique mesur\'e}
Un \textbf{syst\`eme dynamique mesur\'e} est un quadruplet
$(X, \mathcal{B}, \mu, T)$ o\`u $(X, \mathcal{B}, \mu)$ est un espace de
probabilit\'e et $T : X \to X$ est une transformation mesurable
\textbf{pr\'eservant la mesure}\index{transformation pr\'eservant la mesure} :
$\mu(T^{-1}(B)) = \mu(B)$ pour tout $B \in \mathcal{B}$.
\end{definition}

\begin{exemple}[Rotation irrationnelle du cercle]
\index{rotation irrationnelle}
Soit $T_\alpha : \R/\Z \to \R/\Z$ d\'efinie par $T_\alpha(x) = x + \alpha \pmod{1}$
avec $\alpha \notin \Q$. La mesure de Lebesgue sur $[0, 1)$ est invariante par
$T_\alpha$.
\end{exemple}

\begin{exemple}[Application de doublement]
L'application $T(x) = 2x \pmod{1}$ sur $[0, 1)$ pr\'eserve la mesure de
Lebesgue. C'est un endomorphisme (non inversible) avec entropie $\ln 2$.
\end{exemple}

\begin{exemple}[Transformation du boulanger]
\index{transformation du boulanger}
La \textbf{transformation du boulanger} sur $[0,1)^2$ est
\[
    B(x, y) = \begin{cases}
    (2x, y/2) & \text{si } 0 \leq x < 1/2, \\
    (2x - 1, (y+1)/2) & \text{si } 1/2 \leq x < 1.
    \end{cases}
\]
Elle pr\'eserve la mesure de Lebesgue et mod\'elise le p\'etrissage :
\'etirement horizontal, compression verticale, puis repliement.
\end{exemple}

\section{Ergodicit\'e}

\begin{definition}[Ensemble invariant]
Un ensemble $A \in \mathcal{B}$ est \textbf{$T$-invariant} si
$T^{-1}(A) = A$ (modulo un ensemble de mesure nulle).
\end{definition}

\begin{definition}[Transformation ergodique]
\index{ergodicit\'e}
$T$ est \textbf{ergodique} si tout ensemble $T$-invariant a mesure $0$ ou $1$.
De mani\`ere \'equivalente, toute fonction mesurable $T$-invariante
($\varphi \circ T = \varphi$ p.p.) est constante p.p.
\end{definition}

\begin{theoreme}[Ergodicit\'e de la rotation irrationnelle]
La rotation $T_\alpha$ est ergodique par rapport \`a la mesure de Lebesgue si
et seulement si $\alpha$ est irrationnel.
\end{theoreme}

\begin{proof}
Soit $\varphi \in L^2([0,1))$ invariante : $\varphi(x + \alpha) = \varphi(x)$ p.p.
En d\'ecomposant en s\'erie de Fourier $\varphi(x) = \sum_{n \in \Z} c_n e^{2\pi i n x}$,
l'invariance donne $c_n e^{2\pi i n \alpha} = c_n$ pour tout $n$. Si $\alpha \notin \Q$,
alors $e^{2\pi i n \alpha} \neq 1$ pour $n \neq 0$, donc $c_n = 0$ pour
tout $n \neq 0$. Ainsi $\varphi = c_0$ est constante p.p.
\end{proof}

\section{Le th\'eor\`eme ergodique de Birkhoff}

\begin{theoreme}[Birkhoff, 1931]
\label{thm:birkhoff}
\index{th\'eor\`eme de Birkhoff}
Soit $(X, \mathcal{B}, \mu, T)$ un syst\`eme dynamique mesur\'e et
$\varphi \in L^1(X, \mu)$. Alors, pour $\mu$-presque tout $x$ :
\[
    \overline{\varphi}(x) \coloneqq \lim_{N \to \infty} \frac{1}{N}
    \sum_{n=0}^{N-1} \varphi(T^n(x))
\]
existe et $\overline{\varphi}$ est $T$-invariante. De plus,
$\int_X \overline{\varphi} \, \mathrm{d}\mu = \int_X \varphi \, \mathrm{d}\mu$.
\end{theoreme}

\begin{corollaire}[Cas ergodique]
Si $T$ est ergodique, alors pour toute $\varphi \in L^1$ et $\mu$-presque
tout $x$ :
\[
    \lim_{N \to \infty} \frac{1}{N} \sum_{n=0}^{N-1} \varphi(T^n(x))
    = \int_X \varphi \, \mathrm{d}\mu.
\]
La \textbf{moyenne temporelle \'egale la moyenne spatiale}.
\end{corollaire}

\begin{formules}[title=\textbf{Th\'eor\`emes ergodiques fondamentaux}]
\begin{itemize}
    \item \textbf{Von Neumann (1932)} : convergence en norme $L^2$ de
    $\frac{1}{N}\sum_{n=0}^{N-1} \varphi \circ T^n$ vers la projection sur
    les fonctions $T$-invariantes.
    \item \textbf{Birkhoff (1931)} : convergence presque s\^ure (plus fort).
    \item \textbf{Kingman (1968)} : th\'eor\`eme ergodique sous-additif pour
    les suites $(a_n)$ avec $a_{m+n} \leq a_m + a_n \circ T^m$.
\end{itemize}
\end{formules}

\begin{exemple}[\'Equir\'epartition de Weyl]
En appliquant Birkhoff \`a $T_\alpha$ avec $\varphi = \ind_{[a,b]}$, on obtient
le th\'eor\`eme de Weyl : pour $\alpha \notin \Q$,
\[
    \lim_{N \to \infty} \frac{\#\{0 \leq n < N : \{n\alpha\} \in [a, b]\}}{N}
    = b - a.
\]
La suite $(n\alpha \bmod 1)_{n \geq 0}$ est \'equir\'epartie modulo 1.
\end{exemple}

\section{M\'elange}

\begin{definition}[M\'elange faible et fort]
\index{m\'elange}
$T$ est \textbf{faiblement m\'elangeante} si pour tous $A, B \in \mathcal{B}$ :
\[
    \lim_{N \to \infty} \frac{1}{N} \sum_{n=0}^{N-1}
    \abs{\mu(T^{-n}(A) \cap B) - \mu(A)\mu(B)} = 0.
\]
$T$ est \textbf{(fortement) m\'elangeante} si :
\[
    \lim_{n \to \infty} \mu(T^{-n}(A) \cap B) = \mu(A)\mu(B).
\]
\end{definition}

\begin{proposition}[Implications]
\[
    \text{m\'elange fort} \implies \text{m\'elange faible} \implies \text{ergodicit\'e}.
\]
Aucune implication inverse n'est vraie en g\'en\'eral.
\end{proposition}

\begin{exemple}
L'application de doublement $T(x) = 2x \pmod{1}$ et la transformation du
boulanger sont fortement m\'elangeantes. La rotation irrationnelle est
ergodique mais \textbf{non} m\'elangeante (m\^eme faiblement).
\end{exemple}

\begin{theoreme}[Caract\'erisation spectrale du m\'elange]
$T$ est fortement m\'elangeante si et seulement si pour tous
$\varphi, \psi \in L^2(X, \mu)$ :
\[
    \lim_{n \to \infty} \int_X (\varphi \circ T^n) \cdot \psi \, \mathrm{d}\mu
    = \int_X \varphi \, \mathrm{d}\mu \cdot \int_X \psi \, \mathrm{d}\mu.
\]
\end{theoreme}

\section{D\'ecomposition ergodique}

\begin{theoreme}[D\'ecomposition ergodique]
\index{d\'ecomposition ergodique}
Soit $T : X \to X$ pr\'eservant $\mu$ sur un espace de probabilit\'e standard.
Alors il existe une d\'ecomposition (essentiellement unique)
\[
    \mu = \int_Y \mu_y \, \mathrm{d}\nu(y)
\]
o\`u $\nu$ est une mesure de probabilit\'e sur un espace mesurable $Y$ et
chaque $\mu_y$ est une mesure de probabilit\'e $T$-invariante et
\textbf{ergodique}.
\end{theoreme}

\begin{remarque}
La d\'ecomposition ergodique signifie que tout syst\`eme dynamique mesur\'e
peut \^etre d\'ecompos\'e en \og composantes ergodiques \fg irr\'eductibles.
C'est l'analogue de la d\'ecomposition d'un espace de Hilbert en sous-espaces
irr\'eductibles.
\end{remarque}

\section{Entropie m\'etrique}

\begin{definition}[Entropie de Kolmogorov-Sinai]
\index{entropie de Kolmogorov-Sinai}
Soit $\mathcal{P} = \{P_1, \ldots, P_k\}$ une partition mesurable de $X$.
L'\textbf{entropie} de $\mathcal{P}$ est $H(\mathcal{P}) = -\sum_{i=1}^k
\mu(P_i) \ln \mu(P_i)$. L'\textbf{entropie m\'etrique} de $T$ par rapport
\`a $\mathcal{P}$ est
\[
    h_\mu(T, \mathcal{P}) = \lim_{n \to \infty} \frac{1}{n}
    H\left(\bigvee_{k=0}^{n-1} T^{-k}\mathcal{P}\right),
\]
et l'\textbf{entropie de Kolmogorov-Sinai} est
$h_\mu(T) = \sup_{\mathcal{P}} h_\mu(T, \mathcal{P})$.
\end{definition}

\begin{theoreme}[Kolmogorov-Sinai]
Si $\mathcal{P}$ est une partition g\'en\'eratrice (i.e.,
$\bigvee_{n=0}^{\infty} T^{-n}\mathcal{P}$ engendre $\mathcal{B}$
modulo les ensembles de mesure nulle), alors
$h_\mu(T) = h_\mu(T, \mathcal{P})$.
\end{theoreme}

\section{Applications \`a la th\'eorie des nombres}

\begin{exemple}[Fractions continues et transformation de Gauss]
\index{transformation de Gauss}
La \textbf{transformation de Gauss} $G : (0, 1] \to [0, 1)$ d\'efinie par
$G(x) = \{1/x\}$ (partie fractionnaire de $1/x$) encode l'algorithme des
fractions continues. La mesure de Gauss
$\mathrm{d}\mu_G = \frac{1}{\ln 2} \frac{\mathrm{d}x}{1 + x}$
est $G$-invariante et ergodique.
\end{exemple}

\begin{corollaire}[Th\'eor\`eme de Khintchine]
Pour presque tout r\'eel $x \in (0,1)$ dont le d\'eveloppement en fraction
continue est $x = [a_1, a_2, \ldots]$ :
\[
    \lim_{n \to \infty} (a_1 a_2 \cdots a_n)^{1/n}
    = \prod_{k=1}^{\infty} \left(1 + \frac{1}{k(k+2)}\right)^{\ln k / \ln 2}
    \approx 2{,}685\ldots
\]
(constante de Khintchine).
\end{corollaire}

\begin{exemple}[Th\'eor\`eme de L\'evy]
Par le th\'eor\`eme ergodique appliqu\'e \`a $G$ :
\[
    \lim_{n \to \infty} \frac{1}{n} \ln q_n = \frac{\pi^2}{12 \ln 2}
    \approx 1{,}1866\ldots
\]
o\`u $q_n$ est le d\'enominateur du $n$-\`eme convergent. Ce r\'esultat
donne le taux d'approximation typique par les fractions continues.
\end{exemple}

\section{Exercices}

\begin{exercice}[Ergodicit\'e]
Montrer que la transformation du boulanger est ergodique en utilisant les
coefficients de Fourier.
\end{exercice}

\begin{exercice}[Th\'eor\`eme de r\'ecurrence de Poincar\'e]
D\'emontrer que si $T$ pr\'eserve $\mu$ et $\mu(A) > 0$, alors pour
$\mu$-presque tout $x \in A$, il existe $n \geq 1$ tel que $T^n(x) \in A$.
\end{exercice}

\begin{exercice}[M\'elange de l'application de doublement]
Montrer que $T(x) = 2x \pmod{1}$ est fortement m\'elangeante en utilisant
la caract\'erisation spectrale avec les fonctions $e^{2\pi i k x}$.
\end{exercice}

\begin{exercice}[Entropie]
Calculer l'entropie m\'etrique de la transformation de Gauss $G(x) = \{1/x\}$
par rapport \`a la mesure de Gauss. (R\'eponse : $h_{\mu_G}(G) = \pi^2/(6\ln 2)$.)
\end{exercice}

\begin{exercice}[\'Equir\'epartition]
Utiliser le th\'eor\`eme de Birkhoff pour montrer que la suite
$(\sin(n) \bmod 1)_{n \geq 1}$ n'est pas n\'ecessairement \'equir\'epartie
et discuter les conditions sur $\alpha$ pour que $(n^\alpha \bmod 1)$ le soit.
\end{exercice}

\begin{exercice}[Transformation de Gauss]
V\'erifier num\'eriquement la constante de Khintchine en calculant la moyenne
g\'eom\'etrique des coefficients de la fraction continue de nombres al\'eatoires.
\end{exercice}
